\chapter{Twierdzenie o funkcji odwrotnej}
\label{ma2:ch:inverse-function}

Niech $E,F$ będą rzeczywistymi przestrzeniami Banacha, a
$U\subseteq E$ zbiorem otwartym. Badamy, kiedy odwzorowanie
$f:U\to F$ można lokalnie odwrócić i jak obliczać pochodne
jego odwrotności. Wykorzystamy twierdzenia rozdziałów 07,
13--15 oraz wyniki udowodnione w analizie funkcjonalnej.
Odwracalność pochodnej w punkcie nie oznacza jeszcze
odwracalności odwzorowania na całej dziedzinie.

\section{Lokalna odwracalność}

\begin{definicja}[Dyfeomorfizm i lokalny dyfeomorfizm]
\label{ma2:def:local-diffeomorphism}
Niech $U\subseteq E$ i $V\subseteq F$ będą otwarte,
a $m\in\N$, $m\ge1$, lub $m=\infty$.
Dyfeomorfizmem klasy $C^m$ nazywamy bijekcję
$f:U\to V$ taką, że $f\in C^m(U,F)$ oraz
$f^{-1}\in C^m(V,E)$.
Mówimy, że $f:U\to F$ jest lokalnym dyfeomorfizmem
klasy $C^m$ w punkcie $p\in U$, jeżeli istnieje otwarte
otoczenie $W\subseteq U$ punktu $p$ takie, że
$f(W)$ jest otwarty w $F$, a $f|_W:W\to f(W)$
jest dyfeomorfizmem klasy $C^m$.
Jeżeli warunek ten zachodzi w każdym punkcie $U$,
mówimy o lokalnym dyfeomorfizmie na $U$.
\end{definicja}

\begin{twierdzenie}[Twierdzenie o funkcji odwrotnej]
\label{ma2:thm:inverse-function}
Niech $f\in C^1(U,F)$, $p\in U$ oraz niech
$A=Df(p):E\to F$ będzie izomorfizmem liniowym
i topologicznym. Istnieją wtedy otwarte otoczenia
$U_0\subseteq U$ punktu $p$ oraz $V_0\subseteq F$
punktu $q=f(p)$ takie, że
$f|_{U_0}:U_0\to V_0$ jest dyfeomorfizmem klasy $C^1$.
Jeżeli $g:V_0\to U_0$ jest jego odwrotnością,
to dla każdego $y\in V_0$ operator $Df(g(y))$
jest odwracalny oraz
\[
 Dg(y)=(Df(g(y)))^{-1}.
\]
Jeżeli ponadto $f$ jest klasy $C^m$ dla $m\ge2$
lub klasy $C^\infty$, to $g$ jest tej samej klasy.
\end{twierdzenie}

Pełny dowód pierwszej części znajduje się w
\atlasref{FA-15::thm:inverse15}{FA-15: twierdzenie o funkcji odwrotnej},
a dowód regularności wyższych rzędów w
\atlasref{FA-15::cor:inverse-ck15}{FA-15: wersja klasy $C^m$}.
Hipotezy obu wyników są dokładnie hipotezami przyjętymi
powyżej. W dowodzie sprowadza się równanie $f(x)=y$
do równania punktu stałego dla odwzorowania
\[
 T_y(x)=x-A^{-1}(f(x)-y).
\]
Ciągłość $Df$ w normie operatorowej zapewnia, że na
dostatecznie małej domkniętej kuli wokół $p$
odwzorowanie $T_y$ jest kontrakcją. Zupełność
tej kuli umożliwia zastosowanie zasady kontrakcji.
W kolejnym stwierdzeniu wyodrębnimy oszacowania,
które pozwalają ilościowo opisać ten mechanizm.

\begin{stwierdzenie}[Oszacowania lokalne]
\label{ma2:prop:inverse-local-estimates}
Niech $f\in C^1(U,F)$, $E\ne\{0\}$, $p\in U$,
$A=Df(p)$ będzie izomorfizmem liniowym i topologicznym,
$0<\theta<1$ oraz $r>0$. Załóżmy, że
$\overline B(p,r)\subseteq U$ i
\[
 \|I-A^{-1}Df(x)\|\le\theta
 \qquad(x\in\overline B(p,r)).
\]
Wówczas dla $x,z\in\overline B(p,r)$ zachodzi
\[
 \begin{aligned}
 (1-\theta)\|x-z\|&\le
 \|A^{-1}(f(x)-f(z))\|\\
 &\le(1+\theta)\|x-z\|.
 \end{aligned}
\]
Odwzorowanie $f$ jest różnowartościowe na tej kuli,
a jego obraz kuli otwartej zawiera zbiór
\[
 V_r=\{y\in F:\|A^{-1}(y-f(p))\|<(1-\theta)r\}.
\]
Ponadto $f(B(p,r))$ jest otwarty, a odwrotność
$g:f(B(p,r))\to B(p,r)$ jest klasy $C^1$ i spełnia
\[
 \|g(y)-g(w)\|\le
 \frac{\|A^{-1}\|}{1-\theta}\|y-w\|,\qquad
 \|Dg(y)\|\le\frac{\|A^{-1}\|}{1-\theta}.
\]
\end{stwierdzenie}
\begin{proof}
Dla $S(x)=x-A^{-1}f(x)$ mamy
$DS(x)=I-A^{-1}Df(x)$. Odcinek łączący dowolne
$x,z$ z kuli domkniętej pozostaje w tej kuli.
Nierówność o wartości średniej
(twierdzenie \ref{ma2:thm:mean-value-inequality})
daje więc
\[
 \|(x-z)-A^{-1}(f(x)-f(z))\|
 =\|S(x)-S(z)\|\le\theta\|x-z\|.
\]
Stosując nierówność trójkąta w obu kierunkach,
otrzymujemy oba oszacowania normy.
Jeżeli $f(x)=f(z)$, pierwsze oszacowanie daje $x=z$,
ponieważ $1-\theta>0$.

Dla $y\in V_r$ odwzorowanie $T_y$ z dowodu
twierdzenia \ref{ma2:thm:inverse-function} ma na
kuli stałą Lipschitza co najwyżej $\theta$, ponieważ
$T_y(x)-T_y(z)=S(x)-S(z)$. Ponadto
\[
 \|T_y(x)-p\|\le\theta\|x-p\|
   +\|A^{-1}(y-f(p))\|<r
 \quad(x\in\overline B(p,r)).
\]
Zasada kontrakcji, zastosowana tak jak w przywołanym
dowodzie FA-15, daje jedyny punkt stały $x$ w
kuli domkniętej. Z powyższej ostrej nierówności
$x=T_y(x)$ wynika $x\in B(p,r)$.
Ponieważ $A^{-1}$ jest różnowartościowy,
równanie punktu stałego jest równoważne $f(x)=y$.
Stąd $V_r\subseteq f(B(p,r))$.

Dla każdego $x\in B(p,r)$ operator $A^{-1}Df(x)$
różni się od identyczności o operator normy
co najwyżej $\theta<1$. Z twierdzenia
\ref{ma2:thm:inverse-perturbation} zastosowanego
w algebrze $\mathcal L(E)$ wynika jego odwracalność
oraz oszacowanie normy odwrotności przez $1/(1-\theta)$.
Ponieważ $Df(x)=A(A^{-1}Df(x))$,
\[
 (Df(x))^{-1}=(A^{-1}Df(x))^{-1}A^{-1}.
\]
Twierdzenie \ref{ma2:thm:inverse-function},
zastosowane w każdym $x\in B(p,r)$, daje otwarte
otoczenia $x$, zawarte w tej kuli, o otwartych
obrazach. Suma ich obrazów jest całym $f(B(p,r))$,
więc ten obraz jest otwarty. Z różnowartościowości
lokalne odwrotności są ograniczeniami jednej
odwrotności $g$, zatem $g$ jest klasy $C^1$.
Podstawiając $x=g(y)$, $z=g(w)$ do dolnego
oszacowania, otrzymujemy
\[
 (1-\theta)\|g(y)-g(w)\|
 \le\|A^{-1}(y-w)\|\le\|A^{-1}\|\|y-w\|.
\]
Oszacowanie $Dg(y)$ wynika ze wzoru na pochodną
odwrotności i oszacowania normy $(Df(g(y)))^{-1}$.
\end{proof}

Zbiór $V_r$ jest otwarty jako przeciwobraz kuli
otwartej przez ciągłe odwzorowanie afiniczne.
Zawiera zwykłą kulę w $F$ o środku $f(p)$ i
promieniu $(1-\theta)r/\|A^{-1}\|$, ponieważ
$\|A^{-1}(y-f(p))\|\le\|A^{-1}\|\|y-f(p)\|$.
Z ciągłości $Df$ w $p$ zawsze można dobrać
$r>0$ spełniające hipotezy stwierdzenia dla dowolnego
ustalonego $\theta\in(0,1)$.

\section{Pochodne odwrotności}

\begin{stwierdzenie}[Warunek konieczny odwracalności różniczkowej]
\label{ma2:prop:inverse-necessary}
Niech $U\subseteq E$ i $V\subseteq F$ będą otwarte,
a $f:U\to V$ bijekcją o odwrotności $g:V\to U$.
Jeżeli $f$ jest różniczkowalna w $p\in U$,
a $g$ w $q=f(p)$, to $Df(p)$ jest izomorfizmem
liniowym i topologicznym oraz $Dg(q)=Df(p)^{-1}$.
\end{stwierdzenie}
\begin{proof}
Z tożsamości $g\circ f=\id_U$ i $f\circ g=\id_V$
oraz reguły łańcuchowej
(twierdzenie \ref{ma2:thm:frechet-chain}) wynika
\[
 Dg(q)Df(p)=I_E,\qquad Df(p)Dg(q)=I_F.
\]
Pochodne Frécheta są ograniczonymi operatorami
liniowymi. Pierwsza równość daje lewą odwrotność,
a druga prawą odwrotność, więc $Df(p)$ jest
bijekcją o ograniczonej odwrotności $Dg(q)$.
\end{proof}

\begin{stwierdzenie}[Druga pochodna odwrotności]
\label{ma2:prop:inverse-second-derivative}
Niech $f:U_0\to V_0$ będzie dyfeomorfizmem klasy
$C^2$ o odwrotności $g$. Dla $y\in V_0$ oznaczmy
$x=g(y)$ i $B=(Df(x))^{-1}$. Wtedy dla $u,v\in F$
\[
 D^2g(y)[u,v]=-B\bigl(D^2f(x)[Bu,Bv]\bigr).
\]
W szczególności
$\|D^2g(y)\|\le\|B\|^3\|D^2f(x)\|$.
\end{stwierdzenie}
\begin{proof}
Różniczkujemy dwukrotnie $f\circ g=\id_{V_0}$
korzystając ze stwierdzenia \ref{ma2:prop:second-chain}.
Druga pochodna identyczności jest zerowa, zatem
\[
 0=D^2f(x)[Dg(y)u,Dg(y)v]
           +Df(x)D^2g(y)[u,v].
\]
Ze stwierdzenia \ref{ma2:prop:inverse-necessary}
mamy $Dg(y)=B$. Stosując operator $B$ do obu stron
równości, otrzymujemy żądany wzór.
Dalej z normy operatorowej i wieloliniowej
\[
 \|D^2g(y)[u,v]\|
 \le\|B\|\|D^2f(x)\|\|Bu\|\|Bv\|
 \le\|B\|^3\|D^2f(x)\|\|u\|\|v\|.
\]
Bierzemy supremum po $\|u\|,\|v\|\le1$.
\end{proof}

\section{Przypadek euklidesowy}

\begin{wniosek}[Kryterium wyznacznika]
\label{ma2:cor:inverse-jacobian}
Niech $n\ge1$, $U\subseteq\R^n$ będzie otwarty,
$f\in C^1(U,\R^n)$ i $p\in U$. Wówczas $f$ jest
lokalnym dyfeomorfizmem klasy $C^1$ w $p$ wtedy
i tylko wtedy, gdy $\det J_f(p)\ne0$, gdzie
$J_f(p)=(\partial_j f_i(p))_{i,j=1}^n$.
Dla lokalnej odwrotności $g$ zachodzi
\[
 J_g(y)=J_f(g(y))^{-1},\qquad
 \det J_g(y)=\frac{1}{\det J_f(g(y))}.
\]
\end{wniosek}
\begin{proof}
Macierz $J_f(p)$ reprezentuje operator $Df(p)$
w bazach standardowych. Wyznacznik jest niezerowy
dokładnie wtedy, gdy macierz jest odwracalna.
W skończonym wymiarze odwrotna macierz reprezentuje
ograniczony operator liniowy. Warunek dostateczny
wynika więc z twierdzenia \ref{ma2:thm:inverse-function},
a konieczny ze stwierdzenia
\ref{ma2:prop:inverse-necessary}.
Macierz pierwszej pochodnej odwrotności jest macierzą
odwrotną do $J_f(g(y))$. Biorąc wyznacznik równości
$J_g(y)J_f(g(y))=I_n$, otrzymujemy drugą formułę.
\end{proof}

\begin{przyklad}[Odwrotność bez jawnego wzoru]
\label{ma2:ex:inverse-polynomial}
Niech $f:\R^2\to\R^2$ będzie dane wzorem
$f(x_1,x_2)=(x_1+x_2^2,x_2+x_1^2)$.
Mamy
\[
 J_f(x_1,x_2)=
 \begin{pmatrix}1&2x_2\\2x_1&1\end{pmatrix},
 \qquad\det J_f(x_1,x_2)=1-4x_1x_2.
\]
W początku $f(0)=0$ i $J_f(0)=I_2$, więc istnieje
lokalna odwrotność $g$ klasy $C^\infty$, z $Dg(0)=I_2$.
Różniczkując obie składowe $f$ ponownie, otrzymujemy
$D^2f(0)[u,v]=(2u_2v_2,2u_1v_1)$.
Ze stwierdzenia \ref{ma2:prop:inverse-second-derivative}
wynika $D^2g(0)[u,v]=(-2u_2v_2,-2u_1v_1)$.
Wzór Taylora z resztą Peana
(twierdzenie \ref{ma2:thm:taylor-peano}) daje
\[
 g(y_1,y_2)=(y_1-y_2^2,y_2-y_1^2)+o(\|y\|^2)
 \qquad(y\to0).
\]
Możemy więc wyznaczyć lokalne przybliżenie odwrotności
bez rozwiązywania układu nieliniowego w postaci jawnej.
\end{przyklad}

\begin{przyklad}[Współrzędne lokalne]
\label{ma2:ex:inverse-shear}
Dla $f(x,y)=(x,y+x^2)$ mamy
\[
 J_f(x,y)=\begin{pmatrix}1&0\\2x&1\end{pmatrix},
 \qquad\det J_f(x,y)=1.
\]
W tym przypadku odwrotność istnieje globalnie:
$g(u,v)=(u,v-u^2)$, ponieważ bezpośrednio
$f(g(u,v))=(u,v)$ oraz $g(f(x,y))=(x,y)$.
Proste $x=a$ przechodzą na proste $u=a$,
natomiast proste $y=b$ na parabole $v=b+u^2$.
Rysunek \ref{ma2:fig:inverse-shear} ilustruje
obraz kwadratu $[-1,1]^2$.
\end{przyklad}

\begin{figure}[htbp]
\centering
\begin{tikzpicture}[scale=1.65,>=Stealth]
\begin{scope}[xshift=-1.85cm]
\draw[->,gray] (-1.25,0)--(1.3,0) node[right] {$x$};
\draw[->,gray] (0,-1.25)--(0,1.3) node[above] {$y$};
\foreach \a in {-0.5,0,0.5}{
  \draw[blue!50,thin] (\a,-1)--(\a,1);
  \draw[blue!50,thin] (-1,\a)--(1,\a);
}
\draw[blue!70!black,thick] (-1,-1) rectangle (1,1);
\end{scope}
\draw[->,thick] (-0.48,0.45)--(0.48,0.45)
 node[midway,above] {$f$};
\begin{scope}[xshift=1.85cm]
\draw[->,gray] (-1.25,0)--(1.3,0) node[right] {$u$};
\draw[->,gray] (0,-1.25)--(0,2.25) node[above] {$v$};
\foreach \a in {-0.5,0,0.5}{
  \draw[blue!50,thin] (\a,{\a*\a-1})--(\a,{\a*\a+1});
  \draw[blue!50,thin] plot[domain=-1:1,samples=35,variable=\t]
      (\t,{\t*\t+\a});
}
\draw[blue!70!black,thick] (-1,0)--(-1,2);
\draw[blue!70!black,thick] (1,0)--(1,2);
\draw[blue!70!black,thick] plot[domain=-1:1,samples=35,variable=\t]
    (\t,{\t*\t-1});
\draw[blue!70!black,thick] plot[domain=-1:1,samples=35,variable=\t]
    (\t,{\t*\t+1});
\end{scope}
\end{tikzpicture}
\caption{Obraz siatki współrzędnych przez $f(x,y)=(x,y+x^2)$.}
\label{ma2:fig:inverse-shear}
\end{figure}

\section{Odwracalność lokalna a globalna}

\begin{wniosek}[Sklejanie lokalnych odwrotności]
\label{ma2:cor:inverse-global-injective}
Niech $f\in C^m(U,F)$, gdzie $1\le m\le\infty$,
i niech $Df(x)$ będzie izomorfizmem liniowym
i topologicznym dla każdego $x\in U$.
Wtedy $f$ jest odwzorowaniem otwartym na $U$.
Jeżeli dodatkowo $f$ jest różnowartościowe, to
$f:U\to f(U)$ jest dyfeomorfizmem klasy $C^m$.
\end{wniosek}
\begin{proof}
Niech $W\subseteq U$ będzie otwarty. Dla dowolnego
$y\in f(W)$ wybieramy $x\in W$ z $f(x)=y$.
Stosujemy twierdzenie \ref{ma2:thm:inverse-function}
do ograniczenia $f|_W$ w punkcie $x$.
Otrzymujemy otwarte otoczenie $y$ zawarte w $f(W)$.
Zatem $f(W)$ jest otwarty.
Jeżeli $f$ jest różnowartościowe, to $f:U\to f(U)$
jest bijekcją. Każda lokalna odwrotność zgadza się
na swojej dziedzinie z odwrotnością globalną:
dla $y$ z tej dziedziny oba obrazy odwrotne są
rozwiązaniami równania $f(x)=y$, a rozwiązanie
jest jedyne. W konsekwencji globalna odwrotność
jest lokalnie klasy $C^m$ w każdym punkcie $f(U)$,
czyli należy do $C^m(f(U),E)$.
\end{proof}

\begin{przyklad}[Niezerowy wyznacznik nie daje globalnej różnowartościowości]
\label{ma2:ex:inverse-polar-exponential}
Rozważmy odwzorowanie $f:\R^2\to\R^2$ dane przez
\[
 f(s,t)=(e^s\cos t,e^s\sin t).
\]
Jego macierz Jacobiego i wyznacznik są równe
\[
 J_f(s,t)=e^s
 \begin{pmatrix}\cos t&-\sin t\\\sin t&\cos t\end{pmatrix},
 \qquad \det J_f(s,t)=e^{2s}>0.
\]
Na mocy wniosku \ref{ma2:cor:inverse-jacobian}
jest to lokalny dyfeomorfizm klasy $C^\infty$
w każdym punkcie. Jednak $f(s,t+2\pi)=f(s,t)$
dla wszystkich $(s,t)$, więc nie jest różnowartościowe.

Na pasie $U_0=\R\times(-\pi/2,\pi/2)$ obrazem jest
prawa półpłaszczyzna $V_0=\{(u,v)\in\R^2:u>0\}$.
Dla $(s,t)\in U_0$ pierwsza współrzędna jest dodatnia.
Odwrotnie, dla $u>0$ połóżmy
\[
 g(u,v)=
 \left(\tfrac12\log(u^2+v^2),\arctan(v/u)\right).
\]
Jeżeli $r=\sqrt{u^2+v^2}$ i $\alpha=\arctan(v/u)$,
to $\alpha\in(-\pi/2,\pi/2)$,
$\cos\alpha=u/r$, $\sin\alpha=v/r$.
Stąd $f(g(u,v))=(u,v)$.
Z drugiej strony $u^2+v^2=e^{2s}$ i $v/u=\tan t$,
a $\arctan(\tan t)=t$ na rozważanym pasie.
Zatem $g(f(s,t))=(s,t)$ i $g$ jest odwrotnością
ograniczenia. Jej macierz Jacobiego wynosi
\[
 J_g(u,v)=\frac{1}{u^2+v^2}
 \begin{pmatrix}u&v\\-v&u\end{pmatrix}.
\]
Istotnie, pierwszy wiersz otrzymujemy różniczkując
$\tfrac12\log(u^2+v^2)$, drugi z reguły
łańcuchowej i wzoru $(\arctan z)'=1/(1+z^2)$.
Wyznacznik $J_g$ jest równy $1/(u^2+v^2)$,
zgodnie z wnioskiem \ref{ma2:cor:inverse-jacobian}.
\end{przyklad}

\begin{przyklad}[Odwrotność ciągła nie musi być różniczkowalna]
\label{ma2:ex:inverse-cubic}
Funkcja $f:\R\to\R$, $f(x)=x^3$, jest ściśle
rosnąca i jest bijekcją o ciągłej odwrotności
$g(y)=\sqrt[3]{y}$. Jednak $Df(0)=0$.
Dla $y\ne0$ iloraz różnicowy odwrotności w zerze
wynosi
\[
 \frac{g(y)-g(0)}{y}=|y|^{-2/3}\longrightarrow+\infty
 \qquad(y\to0).
\]
Nie istnieje więc skończona pochodna odwrotności
w zerze. Nie przeczy to twierdzeniu
\ref{ma2:thm:inverse-function}, ponieważ jego
hipoteza odwracalności pochodnej nie jest spełniona.
\end{przyklad}

\begin{przyklad}[Pochodna odwracalna w jednym punkcie bez lokalnej iniektywności]
\label{ma2:ex:inverse-discontinuous-derivative}
Niech $f:\R\to\R$ będzie określona wzorem
\[
 f(x)=
 \begin{cases}
 x+2x^2\sin(1/x),&x\ne0,\\
 0,&x=0.
 \end{cases}
\]
Funkcja jest różniczkowalna wszędzie, a w zerze
$f'(0)=1$, ponieważ
$(f(h)-f(0))/h=1+2h\sin(1/h)\to1$.
Dla $x\ne0$ reguły iloczynu i łańcuchowa dają
\[
 f'(x)=1+4x\sin(1/x)-2\cos(1/x).
\]
Dla $a_n=1/(2\pi n)$ oraz $b_n=1/((2n+1)\pi)$,
$n\ge1$, otrzymujemy $f'(a_n)=-1$ oraz $f'(b_n)=3$.
Oba ciągi dążą do zera, więc $f'$ nie jest
ciągła w zerze. Gdyby $f$ było różnowartościowe
w pewnym otoczeniu zera, jego ograniczenie do
zawartego w nim otwartego przedziału wokół zera
byłoby ciągłe i różnowartościowe, a więc ściśle
monotoniczne. Wynika to z własności wartości
pośredniej: dla dowolnych trzech punktów
$x_1<x_2<x_3$ wartość $f(x_2)$ musi leżeć ściśle
między $f(x_1)$ i $f(x_3)$, gdyż inaczej pewna
wartość byłaby osiągana na obu sąsiednich odcinkach.
Warunek ten wymusza ten sam kierunek nierówności
dla wszystkich par uporządkowanych punktów.
Pochodna funkcji rosnącej jest nieujemna, a malejącej
niedodatnia, o ile istnieje, co wynika bezpośrednio
ze znaków ilorazów różnicowych.
Tymczasem na każdym takim przedziale znajdują się
odpowiednie $a_n,b_n$, w których pochodne mają
przeciwne znaki. Jest to sprzeczność.
Samo istnienie odwracalnej pochodnej w punkcie
nie wystarcza zatem do tezy twierdzenia o funkcji
odwrotnej.
\end{przyklad}

\begin{przyklad}[Przykład w przestrzeni Hilberta]
\label{ma2:ex:inverse-hilbert}
Niech $H$ będzie rzeczywistą przestrzenią Hilberta,
$v\in H\setminus\{0\}$, a $f:H\to H$ będzie dane przez
$f(x)=x+\|x\|^2v$.
Jest to przykład z
\atlasref{FA-15::ex:inverse-hilbert15}{FA-15: lokalna odwrotność w przestrzeni Hilberta}.
Rozwijając kwadrat normy, otrzymujemy
\[
 f(x+h)-f(x)=h+2\langle x,h\rangle v+\|h\|^2v .
\]
Pierwsze dwa składniki są liniowe i ograniczone
względem $h$, a norma ostatniego po podzieleniu
przez $\|h\|$ wynosi $\|h\|\|v\|\to0$.
Stąd $Df(x)h=h+2\langle x,h\rangle v$.
Pochodna zależy w sposób ciągły i afiniczny od $x$,
więc $f$ jest klasy $C^\infty$. W zerze $Df(0)=I$,
a
\[
 \|I-Df(x)\|\le2\|x\|\|v\|
\]
z nierówności Cauchy'ego--Schwarza.
Dla $r=1/(4\|v\|)$ i $\theta=1/2$ stosujemy
stwierdzenie \ref{ma2:prop:inverse-local-estimates}
z $p=0$, $A=I$. Otrzymujemy różnowartościowość
na $\overline B(0,r)$ oraz istnienie gładkiej
odwrotności na $f(B(0,r))$, który zawiera
$B(0,r/2)$. Stała Lipschitza odwrotności jest
nie większa niż $2$.
Mimo to $f$ nie jest globalnie różnowartościowe:
dla $w=-v/\|v\|^2$ mamy $w\ne0$ i
\[
 f(w)=-v/\|v\|^2+\|v\|^{-2}v=0=f(0).
\]
\end{przyklad}

W następnym rozdziale zastosujemy lokalną
odwracalność do rozwiązywania równań z parametrami.
